\documentclass{article}

\usepackage[utf8]{inputenc}
\usepackage[T1]{fontenc}
\usepackage{helvet}
\usepackage{geometry}
\usepackage{xcolor}
\usepackage{eso-pic}
\usepackage{fancyhdr}
\usepackage{parskip}
\usepackage{microtype}
\usepackage[english, ukrainian]{babel}
\usepackage{url}
\usepackage{amsmath}
\usepackage{amssymb}
\usepackage{amsthm}
\usepackage{longtable}
\usepackage{booktabs}
\usepackage{hyperref}
\usepackage{titlesec}

\definecolor{nato}{HTML}{293757}
\definecolor{footergray}{gray}{0.65}

\geometry{a4paper,left=3cm,right=3cm,top=3cm,bottom=3cm}
\renewcommand{\familydefault}{\sfdefault}
\setcounter{secnumdepth}{3}

\titleformat{\section}
  {\color{nato}\Huge\bfseries}{\thesection.}{1em}{}
\titlespacing*{\section}{0pt}{30pt}{12pt}

\titleformat{\subsection}
  {\color{nato}\LARGE\bfseries}{\thesubsection.}{1em}{}
\titlespacing*{\subsection}{0pt}{20pt}{8pt}

\titleformat{\subsubsection}
  {\color{nato}\Large\bfseries}{\thesubsubsection.}{1em}{}
\titlespacing*{\subsubsection}{0pt}{10pt}{6pt}

\fancyhf{}
\fancyhead[R]{\small\color{footergray} 27 червня 2026}
\fancyfoot[L]{\small\color{footergray} Техноробочий проект ERP/1. Модуль Облік}
\fancyfoot[R]{\small\color{footergray}\thepage}

\newcommand\HUGE[1]{\fontsize{40}{40}\selectfont #1}

\renewcommand{\headrulewidth}{0pt}
\renewcommand{\footrulewidth}{0.4pt}
\renewcommand{\footrule}{\color{footergray}\hrule width \textwidth height 0.4pt}

\begin{document}

\pagestyle{fancy}

\begin{titlepage}
  \AddToShipoutPictureBG*{\AtPageLowerLeft{\color{nato}\rule{\paperwidth}{\paperheight}}}
  \color{white}\thispagestyle{empty}\vspace*{3cm}\centering
  {\Huge  \textbf{ERP/1: Облік} \par} \vspace{1.5cm}
  {\HUGE  \textbf{Техноробочий проект} \par} \vspace{1.5cm}
  {\LARGE \textbf{на розробку та впровадження \\ підсистеми обліку, кадрів \\  та заробітної плати} \par} \vspace{1.5cm}
  {\large Відповідно до ДСТУ 34.201-89 та Постанови КМУ № 205 від 21.02.2025 \par}
  {\large Формалізація вимог за стандартом ISO/IEC/IEEE 29148:2018 \par}
  \vfill
  {\large 2026 \copyright\ Системи Електронного Урядування \par}
\end{titlepage}

\newpage

\begin{center}
  {\Huge\color{nato}\bfseries Зміст\par}
\end{center}
\vspace{40pt}
\tableofcontents

\newpage

\section{Загальні відомості про засіб інформатизації}

\subsection{Найменування та підстави для розробки}
\textbf{Найменування засобу інформатизації:} Підсистема «ERP/1: Облік» (далі --- Модуль «Облік»).
\\
\textbf{Відомості про замовника:} ТОВ "Системи Електронного Урядування" в ТОВ "Криптографічні телесистеми".
\\
\textbf{Підстави для виконання робіт:}
\begin{itemize}
    \item Бюджетний кодекс України;
    \item Закон України «Про бухгалтерський облік та фінансову звітність в Україні»;
    \item Постанова Кабінету Міністрів України від 21 лютого 2025 року № 205 «Деякі питання створення, адміністрування та забезпечення функціонування засобу інформатизації».
\end{itemize}

\subsection{Призначення та цілі створення засобу}
Модуль «Облік» призначений для комплексної автоматизації фінансової та господарської діяльності з можливістю багатовимірного аналітичного обліку, розрахунку заробітної плати та складання регламентованих звітів.
\\
\textbf{Цілі створення засобу:}
\begin{itemize}
    \item Автоматизація ведення Головної книги та бухгалтерських проводок;
    \item Складання оборотного балансу та фінансових звітів у реальному часі;
    \item Розрахунок заробітної плати та податків з особових рахунків працівників;
    \item Управління матеріально-технічним забезпеченням, закупівлями та складом ТМЦ;
    \item Контроль виконання кошторисів за кодами КЕКВ;
    \item Забезпечення автентифікації та підписання відомостей за допомогою КЕП (CAdES-X Long).
\end{itemize}

\subsection{План-графік виконання етапів робіт}
\begin{enumerate}
    \item \textbf{Етап 1: Проектування та погодження архітектури} --- розробка та погодження UML/Figma специфікацій. (1 місяць).
    \item \textbf{Етап 2: Створення схем даних} --- реалізація Erlang records, створення таблиць Mnesia та RocksDB сховищ. (1.5 місяці).
    \item \textbf{Етап 3: Розробка бізнес-процесів} --- реалізація BPE DSL для ЗП, проводок та закупівель. (2 місяці).
    \item \textbf{Етап 4: Інтеграційні сервіси} --- налаштування обміну з Є-Казна, ДПС, ПФУ. (1.5 місяці).
    \item \textbf{Етап 5: Верифікація та тестування} --- Unit, Integration тестування, UAT приймання. (1 місяць).
\end{enumerate}

\newpage
\section{Відомості про робочий процес та умови експлуатації}

\subsection{Опис бізнес-процесів (BPMN / FSM)}

\subsubsection{Процес проводки первинного документа}
\begin{enumerate}
    \item \textbf{StartPosting:} Створення проекту документа бухгалтером.
    \item \textbf{EnterData:} Введення реквізитів (сума, рахунки дебету/кредиту, код КЕКВ).
    \item \textbf{ValidateDoubleEntry:} Автоматична перевірка валідатором подвійного запису. У разі успіху --- перехід до \texttt{CheckBudget}, інакше --- повернення до \texttt{EnterData}.
    \item \textbf{CheckBudget:} Автоматичний контроль відповідності лімітам видатків за КЕКВ.
    \item \textbf{SignDoc:} Накладання КЕП бухгалтера на первинний документ.
    \item \textbf{PostToLedger:} Автоматичне проведення суми по рахунках Головної книги.
    \item \textbf{PostingCompleted:} Фіксація проводки.
\end{enumerate}

\subsubsection{Процес розрахунку заробітної плати}
\begin{enumerate}
    \item \textbf{StartPayroll:} Запуск розрахункового періоду за місяць.
    \item \textbf{CollectTimesheets:} Завантаження та затвердження табелів робочого часу працівників.
    \item \textbf{CalculateTaxes:} Автоматичний розрахунок нарахувань, утримань ПДФО, військового збору та ЄСВ.
    \item \textbf{ApproveRegister:} Погодження відомості розрахунку Головним бухгалтером.
    \item \textbf{QESSign:} Накладання КЕП керівника закладу на відомість виплати.
    \item \textbf{DisburseFunds:} Формування платіжного реєстру для відправки в Є-Казна / банк.
    \item \textbf{PayrollCompleted:} Закриття періоду.
\end{enumerate}

\subsubsection{Процес планування закупівель}
\begin{enumerate}
    \item \textbf{StartProcurement:} Запуск процедури планування закупівлі ТМЦ.
    \item \textbf{DraftRequest:} Створення заявки на закупівлю із зазначенням кодів ТМЦ та КЕКВ.
    \item \textbf{VerifyFunds:} Автоматична перевірка наявності кошторисних асигнувань.
    \item \textbf{ApproveByChief:} Погодження закупівлі керівництвом.
    \item \textbf{GenerateContract:} Автоматичне формування договору в реєстрі укладених угод.
    \item \textbf{ProcurementCompleted:} Завершення етапу планування.
\end{enumerate}

\subsubsection{Процес підготовки та подання звітності}
\begin{enumerate}
    \item \textbf{StartReporting:} Запуск формування звітного періоду (місяць, квартал, рік).
    \item \textbf{RunConstructor:} Збирання даних з Головної книги та кадрового обліку.
    \item \textbf{ValidateData:} Контроль взаємозв'язку показників форм звітності.
    \item \textbf{QESSignReport:} Накладання КЕП посадових особ.
    \item \textbf{SubmitToPortal:} Відправка звіту через інтеграційне API до казначейського або податкового порталу.
    \item \textbf{ReportAccepted:} Отримання квитанції про успішне прийняття звіту.
\end{enumerate}

\subsection{Опис ролей та прав доступу (ABAC)}
\begin{itemize}
    \item \textbf{Правило 1 (Бухгалтер):}
    $$\text{subject.roles contains 'Accountant'} \land \text{object.entity\_type == 'PrimaryDocument'}$$
    $$\land\ \text{object.department\_id == subject.department\_id} \land \text{object.status == 'draft'}$$
    $$\rightarrow \text{Permit: read, write, sign.}$$
    \item \textbf{Правило 2 (Кадровик):}
    $$\text{subject.roles contains 'HR'} \land (\text{object.entity\_type == 'EmployeeRecord'} \lor \text{object.entity\_type == 'TimesheetEntry'})$$
    $$\rightarrow \text{Permit: read, write, sign. Deny: delete.}$$
    \item \textbf{Правило 3 (Комірник):}
    $$\text{subject.roles contains 'Storekeeper'} \land \text{object.entity\_type == 'WarehouseItem'}$$
    $$\rightarrow \text{Permit: read, write (тільки прибуткові/видаткові ордери).}$$
    \item \textbf{Правило 4 (Контроль подвійного запису):}
    $$\text{object.entity\_type == 'JournalPosting'} \land \text{debit\_sum == credit\_sum}$$
    $$\rightarrow \text{Permit: post.}$$
    \item \textbf{Правило 5 (Кошторисний контроль):}
    $$\text{object.entity\_type == 'PrimaryDocument'} \land \text{object.amount} \le \text{budget.available\_limit}$$
    $$\rightarrow \text{Permit: approve.}$$
\end{itemize}

\subsection{Умови експлуатації}
\begin{itemize}
    \item Середовище виконання: Erlang/OTP версії 26 та вище.
    \item База даних: Erlang Mnesia як єдине native транзакційне сховище. Використання реляційних SQL баз даних не допускається.
    \item Сховище медіа-об'єктів (первинних скан-копій): S3-сумісне сховище або RocksDB через бібліотеку KVS.
\end{itemize}

\newpage
\section{Інформаційне та програмне забезпечення}

\subsection{Інформаційне забезпечення (Реквізити)}
Усі об'єкти описуються Erlang-рекордами:
\begin{verbatim}
-record(primary_document, {
    id, date, document_number, type, counterparty,
    debit_account, credit_account, amount, vat, status = draft
}).

-record(journal_posting, {
    id, date, document_number, description, debit_account,
    credit_account, amount, kekv, department, status = draft
}).

-record(ledger_account, {
    id, code, name, opening_balance, debit_turnover,
    credit_turnover, closing_balance
}).

-record(employee_record, {
    id, full_name, personnel_number, department, position,
    hire_date, salary, status = active
}).

-record(timesheet_entry, {
    id, employee_name, department, days_worked, days_absent,
    overtime_hours, total_hours, status = draft
}).

-record(payroll_line, {
    id, employee_name, department, position, gross_pay,
    pit, military_tax, pension_fund, other_deductions,
    net_pay, status = calculated
}).

-record(supply_contract, {
    id, contract_number, date, supplier, total_value,
    executed_amount, stages_completed, total_stages,
    end_date, status = draft, kekv
}).

-record(warehouse_item, {
    id, code, name, unit, quantity, price, total_value,
    min_stock, warehouse, stock_level = normal
}).
\end{verbatim}

\paragraph{Повнотекстові специфікації реквізитів:} наведено в додатку А до Технічних вимог (див. [acc.tex](file:///Users/tonpa/depot/erpuno/erp.uno/acc/acc.tex)).

\subsection{Програмне забезпечення (Реалізація)}

\subsubsection{bpe\_bookkeeping\_posting.erl (Процес проведення документа)}
\begin{verbatim}
-module(bpe_bookkeeping_posting).
-include("bpe.hrl").
-compile(export_all).

def() ->
    #process{
        name = 'Primary Document Posting',
        beginEvent = 'StartPosting',
        endEvent = 'PostingCompleted',
        tasks = [
            #beginEvent{name='StartPosting'},
            #userTask{name='EnterData', module=?MODULE},
            #serviceTask{name='ValidateDoubleEntry', module=?MODULE},
            #serviceTask{name='CheckBudget', module=?MODULE},
            #userTask{name='SignDoc', module=?MODULE},
            #serviceTask{name='PostToLedger', module=?MODULE},
            #endEvent{name='PostingCompleted'}
        ],
        flows = [
            #sequenceFlow{name='1', source='StartPosting', target='EnterData'},
            #sequenceFlow{name='2', source='EnterData', 
                         target='ValidateDoubleEntry'},
            #sequenceFlow{name='3', source='ValidateDoubleEntry', 
                         target='CheckBudget', condition="validation_ok"},
            #sequenceFlow{name='4', source='ValidateDoubleEntry', 
                         target='EnterData', condition="validation_failed"},
            #sequenceFlow{name='5', source='CheckBudget', 
                         target='SignDoc', condition="budget_ok"},
            #sequenceFlow{name='6', source='CheckBudget', 
                         target='EnterData', condition="budget_overlimit"},
            #sequenceFlow{name='7', source='SignDoc', target='PostToLedger'},
            #sequenceFlow{name='8', source='PostToLedger', 
                         target='PostingCompleted'}
        ]
    }.

action({serviceTask, 'ValidateDoubleEntry'}, Proc) ->
    case acc_validator:validate_posting(Proc) of
        ok -> {reply, Proc, "validation_ok"};
        _  -> {reply, Proc, "validation_failed"}
    end;
action(_, Proc) -> {reply, Proc}.
\end{verbatim}

\subsubsection{bpe\_payroll\_calculation.erl (Процес розрахунку ЗП)}
\begin{verbatim}
-module(bpe_payroll_calculation).
-include("bpe.hrl").
-compile(export_all).

def() ->
    #process{
        name = 'Monthly Payroll Calculation',
        beginEvent = 'StartPayroll',
        endEvent = 'PayrollCompleted',
        tasks = [
            #beginEvent{name='StartPayroll'},
            #serviceTask{name='CollectTimesheets', module=?MODULE},
            #serviceTask{name='CalculateTaxes', module=?MODULE},
            #userTask{name='ApproveRegister', module=?MODULE},
            #userTask{name='QESSign', module=?MODULE},
            #serviceTask{name='DisburseFunds', module=?MODULE},
            #endEvent{name='PayrollCompleted'}
        ],
        flows = [
            #sequenceFlow{name='1', source='StartPayroll', 
                         target='CollectTimesheets'},
            #sequenceFlow{name='2', source='CollectTimesheets', 
                         target='CalculateTaxes'},
            #sequenceFlow{name='3', source='CalculateTaxes', 
                         target='ApproveRegister'},
            #sequenceFlow{name='4', source='ApproveRegister', 
                         target='QESSign'},
            #sequenceFlow{name='5', source='QESSign', 
                         target='DisburseFunds'},
            #sequenceFlow{name='6', source='DisburseFunds', 
                         target='PayrollCompleted'}
        ]
    }.
action(_, Proc) -> {reply, Proc}.
\end{verbatim}

\subsubsection{bpe\_procurement\_planning.erl (Процес планування закупівель)}
\begin{verbatim}
-module(bpe_procurement_planning).
-include("bpe.hrl").
-compile(export_all).

def() ->
    #process{
        name = 'Procurement Planning',
        beginEvent = 'StartProcurement',
        endEvent = 'ProcurementCompleted',
        tasks = [
            #beginEvent{name='StartProcurement'},
            #userTask{name='DraftRequest', module=?MODULE},
            #serviceTask{name='VerifyFunds', module=?MODULE},
            #userTask{name='ApproveByChief', module=?MODULE},
            #serviceTask{name='GenerateContract', module=?MODULE},
            #endEvent{name='ProcurementCompleted'}
        ],
        flows = [
            #sequenceFlow{name='1', source='StartProcurement', 
                         target='DraftRequest'},
            #sequenceFlow{name='2', source='DraftRequest', 
                         target='VerifyFunds'},
            #sequenceFlow{name='3', source='VerifyFunds', 
                         target='ApproveByChief', condition="funds_available"},
            #sequenceFlow{name='4', source='VerifyFunds', 
                         target='DraftRequest', condition="insufficient_funds"},
            #sequenceFlow{name='5', source='ApproveByChief', 
                         target='GenerateContract'},
            #sequenceFlow{name='6', source='GenerateContract', 
                         target='ProcurementCompleted'}
        ]
    }.
action(_, Proc) -> {reply, Proc}.
\end{verbatim}

\subsubsection{acc\_validator.erl (Модуль бізнес-валідацій)}
\begin{verbatim}
-module(acc_validator).
-export([validate_posting/1, check_budget_limit/3]).

-record(posting, {debit_amount = 0.0, credit_amount = 0.0}).

validate_posting(Postings) ->
    DebitSum = lists:foldl(fun(P, Acc) -> Acc + P#posting.debit_amount end, 0.0, Postings),
    CreditSum = lists:foldl(fun(P, Acc) -> Acc + P#posting.credit_amount end, 0.0, Postings),
    case DebitSum == CreditSum of
        true -> ok;
        false -> {error, balance_mismatch}
    end.

check_budget_limit(KEKV, Amount, AvailableLimit) ->
    case Amount =< AvailableLimit of
        true -> ok;
        false -> {error, {budget_overlimit, KEKV}}
    end.
\end{verbatim}

\newpage
\section{Вимоги до засобу за ISO/IEC/IEEE 29148}

\subsection{Функціональні вимоги}
\begin{itemize}
    \item \textbf{[REQ-ACC-FUN-001]} Модуль «Облік» повинен забезпечувати ведення плану рахунків, класифікації КЕКВ та довідників ТМЦ.
    \item \textbf{[REQ-ACC-FUN-002]} Модуль «Облік» повинен виконувати автоматичну перевірку балансу дебету та кредиту для кожної господарської проводки.
    \item \textbf{[REQ-ACC-FUN-003]} Модуль «Облік» повинен підтримувати кошторисний контроль --- блокувати проведення первинних документів, якщо сума видатків за кодом КЕКВ перевищує ліміт асигнувань.
    \item \textbf{[REQ-ACC-FUN-004]} Модуль «Облік» повинен забезпечувати автоматичний розрахунок заробітної плати працівників за місяць на основі затверджених табелів обліку робочого часу.
    \item \textbf{[REQ-ACC-FUN-005]} Модуль «Облік» повинен автоматично розраховувати суми утримань ПДФО (18\%), військового збору (1.5\%) та нарахувань ЄСВ (22\%) згідно з чинним законодавством України.
    \item \textbf{[REQ-ACC-FUN-006]} Модуль «Облік» повинен підтримувати підписання розрахункових та платіжних відомостей за допомогою КЕП CAdES-X Long.
    \item \textbf{[REQ-ACC-FUN-007]} Модуль «Облік» повинен забезпечувати автоматичний розрахунок оборотного балансу та відображення залишків по Головній книзі в режимі реального часу.
    \item \textbf{[REQ-ACC-FUN-008]} Модуль «Облік» повинен забезпечувати ведення особових справ співробітників, враховуючи кадрові накази та штатні оклади.
    \item \textbf{[REQ-ACC-FUN-009]} Модуль «Облік» повинен генерувати первинні документи (Акт, Накладна, Розрахунковий листок) у форматах PDF та XLSX згідно з регламентованими формами.
    \item \textbf{[REQ-ACC-FUN-010]} Модуль «Облік» повинен експортувати фінансові звіти до системи «Є-Звітність» за допомогою інтеграційного HTTPS REST API.
\end{itemize}

\subsection{Вимоги до інтерфейсів та дизайну}
\begin{itemize}
    \item \textbf{[REQ-ACC-UI-001]} Користувацький інтерфейс повинен бути сумісним зі стандартом веб-доступності WCAG 2.1 AA.
    \item \textbf{[REQ-ACC-UI-002]} Інтерфейс користувача повинен функціонувати як Single Page Application (SPA), забезпечуючи оновлення фінансових моніторів через Secure WebSockets (WSS).
    \item \textbf{[REQ-ACC-UI-003]} Обмін даними між веб-клієнтом та сервером повинен використовувати архітектурний шаблон Data-on-Wire з передачею об'єктів у форматі JSON.
\end{itemize}

\subsection{Вимоги до безпеки та захисту інформації}
\begin{itemize}
    \item \textbf{[REQ-ACC-SEC-001]} Автентифікація всіх користувачів підсистеми повинна здійснюватися виключно за допомогою кваліфікованого електронного підпису (КЕП) X.509.
    \item \textbf{[REQ-ACC-SEC-002]} Будь-які запити на читання та запис фінансових даних повинні перевірятися підсистемою контролю доступу на основі ABAC-атрибутів суб'єкта та об'єкта.
    \item \textbf{[REQ-ACC-SEC-003]} Система повинна передавати дані виключно через протокол шифрування TLS 1.3 з використанням безпечних криптографічних алгоритмів.
    \item \textbf{[REQ-ACC-SEC-004]} Оскільки реляційні SQL бази даних не використовуються, ризик SQL-ін'єкцій повністю відсутній.
    \item \textbf{[REQ-ACC-SEC-005]} Персональні дані працівників (особові справи, розрахунки ЗП) повинні бути захищені від несанкціонованого доступу згідно із Законом України «Про захист персональних даних».
\end{itemize}

\subsection{Вимоги до продуктивності}
\begin{itemize}
    \item \textbf{[REQ-ACC-PERF-001]} Час розрахунку та відображення оборотного балансу по 50,000 операцій не повинен перевищувати 500 мс.
    \item \textbf{[REQ-ACC-PERF-002]} Пропускна здатність підсистеми Mnesia повинна бути не менше 500 фінансових проводок на секунду на один серверний вузол.
    \item \textbf{[REQ-ACC-PERF-003]} Час відкриття та відображення реєстру документів на 1000 записів на клієнті не повинен перевищувати 1 секунди.
\end{itemize}

\subsection{Вимоги до логування, тестування та супроводу}
\begin{itemize}
    \item \textbf{[REQ-ACC-LOG-001]} Усі операції зміни фінансових станів, коригування проводок та входу користувачів повинні реєструватися у захищеному журналі аудиту (Audit Log).
    \item \textbf{[REQ-ACC-LOG-002]} Логи аудиту повинні передаватися у форматі JSON до Elasticsearch (стек ELK) в режимі реального часу.
    \item \textbf{[REQ-ACC-ADM-001]} Покриття вихідного коду автоматичними Unit-тестами та інтеграційними тестами надійності повинно бути не менше ніж 80\%.
    \item \textbf{[REQ-ACC-ADM-002]} Поставка програмного забезпечення повинна містити повний пакет супровідної документації (Настанова користувача, Настанова адміністратора) згідно з ДСТУ 3008:2015.
\end{itemize}

\end{document}
